\chapter{Final Examination: Academic Year 112}
\label{exam:final-112}

\section{1. Automatic-door controller (100\%)}
\begin{question}
Design a simple automatic-door controller with two inputs, \(D\) and
\(\mathrm{Done}\), and two outputs, \(\mathrm{OP}\) and \(\mathrm{CL}\).
\begin{itemize}
\item \(D\), detect: an infrared detector asserts this signal when a person
is approaching the door.
\item \(\mathrm{Done}\): the motor produces a one-clock-cycle pulse when
the door reaches either the fully open or the fully closed endpoint.
\item \(\mathrm{OP}\), open: command the motor to open the door.
\item \(\mathrm{CL}\), close: command the motor to close the door.
\end{itemize}
The required behavior is:
\begin{enumerate}
\item While the door is closed and \(D=0\), keep it closed.
\item If a person is detected while it is closed, assert \(\mathrm{OP}\).
Keep opening until \(\mathrm{Done}=1\), then deassert \(\mathrm{OP}\).
\item While fully open, start closing after five continuous seconds with
no person detected.
\item Any detection while fully open restarts the five-second interval.
\item If a person is detected while the door is closing, deassert
\(\mathrm{CL}\), assert \(\mathrm{OP}\), and reopen until the fully open
endpoint produces \(\mathrm{Done}=1\).
\item Otherwise, keep closing until \(\mathrm{Done}=1\), then deassert
\(\mathrm{CL}\).
\end{enumerate}
A \(1\,\mathrm{kHz}\) clock is available: one cycle is \(1\,\mathrm{ms}\).
An error of up to \(\pm0.2\,\mathrm{s}\) in the five-second delay is acceptable.
Treat both inputs as synchronous to this clock.

Supply a state diagram, state equations, output equations, and a
gate-level controller drawing. Explain the meaning of every state.
Combinational components such as multiplexers and adders may be used as
clearly labeled blocks. Flip-flops may be drawn without their internal
circuit or clock wires, but their D and Q pins must be labeled.
Sequential components, including counters and timers, must be expanded
into their internal circuit. Hierarchical subcircuits are allowed when
their connections are clear. The original instructions limit the
examination answer to one sheet; this study edition expands the derivation
and diagrams for readability.
\end{question}

\begin{center}
\begin{tikzpicture}[>=Latex]
\node[block,minimum width=35mm,minimum height=22mm] (controller) at (0,0)
{Door\\controller};
\draw[signal] (-3,0.55) node[left] {\(D\)} -- (controller.west |- 0,0.55);
\draw[signal] (-3,-0.55) node[left] {\(\mathrm{Done}\)} -- (controller.west |- 0,-0.55);
\draw[signal] (controller.east |- 0,0.55) -- (3,0.55) node[right] {\(\mathrm{OP}\)};
\draw[signal] (controller.east |- 0,-0.55) -- (3,-0.55) node[right] {\(\mathrm{CL}\)};
\end{tikzpicture}
\end{center}

\section{Solution: synchronous state machine}
All input decisions below are made at a rising clock edge. Outputs are
decoded from the state held after that edge. Thus detecting a person while
closing switches the command to opening at the detecting edge, without an
extra waiting state.

The question does not specify power-up initialization. For a reproducible
implementation, assume the door is initially physically closed and apply
a synchronous initialization signal \(R\) for one clock edge. It clears
the state and timer registers. \(R\) is an implementation initialization
condition, not an additional sensor or a third operational input in the
question.

\begin{center}
\begin{tabular}{llll}
\toprule
State & \(s_1s_0\) & Meaning & \((\mathrm{OP},\mathrm{CL})\)\\
\midrule
\(\mathsf C\) & 00 & Fully closed & \((0,0)\)\\
\(\mathsf U\) & 01 & Opening & \((1,0)\)\\
\(\mathsf H\) & 11 & Fully open; waiting & \((0,0)\)\\
\(\mathsf X\) & 10 & Closing & \((0,1)\)\\
\bottomrule
\end{tabular}
\end{center}
Write \(F=\mathrm{Done}\). Let \(T=1\) mean that the timer contains 4999;
its exact interpretation at the next clock edge is derived below.
An unlabeled input in a transition condition is irrelevant.

\begin{center}
\begin{tikzpicture}[>=Latex,node distance=31mm,
 every state/.style={minimum size=19mm,align=center},font=\small]
\node[state,initial,initial text={\(R\)}] (c) at (0,0) {\(\mathsf C\)\\00};
\node[state] (u) at (5,0) {\(\mathsf U\)\\01};
\node[state] (h) at (5,-3.5) {\(\mathsf H\)\\11};
\node[state] (x) at (0,-3.5) {\(\mathsf X\)\\10};
\path[->]
 (c) edge[loop above] node {\(\overline D\)} (c)
     edge node[above] {\(D\)} (u)
 (u) edge[loop right] node {\(\overline F\)} (u)
     edge node[right] {\(F\)} (h)
 (h) edge[loop below] node {\(D+\overline T\)} (h)
     edge node[below] {\(\overline D T\)} (x)
 (x) edge[loop left] node[left] {\(\overline D\,\overline F\)} (x)
     edge node[left] {\(\overline D F\)} (c)
     edge[bend left=12] node[above,sloped] {\(D\), even if \(F=1\)} (u);
\end{tikzpicture}
\end{center}
The initialization arrow means that \(R=1\) forces \(\mathsf C\) from
every state; this override is omitted from individual arrows.
In \(\mathsf X\), detection takes priority over an endpoint pulse:
\(D=F=1\) requests reopening. In \(\mathsf H\), detection also takes
priority over timeout and restarts the timer.

\section{State and output equations}
First write next-state functions before applying initialization.
The high state bit is one for either a next state of \(\mathsf H\)
or a next state of \(\mathsf X\):
\begin{align*}
f_1
&=s_1s_0+\overline{s_1}s_0F
  +s_1\overline{s_0}\,\overline D\,\overline F.
\end{align*}
The low state bit is one for a next state of \(\mathsf U\) or
\(\mathsf H\). A direct state expansion simplifies to
\begin{align*}
f_0
&=\overline{s_1}\overline{s_0}D+\overline{s_1}s_0
  +s_1s_0(D+\overline T)+s_1\overline{s_0}D\\
&=D+s_0(\overline{s_1}+\overline T).
\end{align*}
Therefore the D inputs of the two state flip-flops are
\[
\boxed{s_1^+=\overline R f_1,\qquad
s_0^+=\overline R f_0.}
\]
The Moore outputs are
\[
\boxed{\mathrm{OP}=\overline{s_1}s_0,\qquad
\mathrm{CL}=s_1\overline{s_0}.}
\]
For every stable state,
\(\mathrm{OP}\,\mathrm{CL}=0\). All four two-bit patterns are assigned
states. The controller therefore needs only two state flip-flops,
in addition to the timer flip-flops.

\begin{center}
\begin{tabular}{llll}
\toprule
Current state & Condition & Next state & Timer action\\
\midrule
\(\mathsf C\) & \(D=0\) / \(D=1\) & \(\mathsf C\) / \(\mathsf U\) & Clear\\
\(\mathsf U\) & \(F=0\) / \(F=1\) & \(\mathsf U\) / \(\mathsf H\) & Clear\\
\(\mathsf H\) & \(D=1\) & \(\mathsf H\) & Clear\\
\(\mathsf H\) & \(D=0,T=0\) & \(\mathsf H\) & Increment\\
\(\mathsf H\) & \(D=0,T=1\) & \(\mathsf X\) & Clear\\
\(\mathsf X\) & \(D=1\) & \(\mathsf U\) & Clear\\
\(\mathsf X\) & \(D=0,F=0\) & \(\mathsf X\) & Clear\\
\(\mathsf X\) & \(D=0,F=1\) & \(\mathsf C\) & Clear\\
\bottomrule
\end{tabular}
\end{center}
In \(\mathsf C\), \(F\) is ignored; in \(\mathsf U\), \(D\) is ignored;
in \(\mathsf H\), \(F\) is ignored. These are intentional consequences
of the required operation, not missing cases.

\clearpage
\section{Gate-level state register and output decoder}
The following diagrams implement the equations directly. Identically
named nets are connected wires. An input bubble denotes inversion; it is
equivalent to an explicit NOT gate on that input. Each DFF symbol has
its D pin on the left and Q pin on the right. All DFFs in this chapter
use the same \(1\,\mathrm{kHz}\) rising-edge clock; the clock wires are
omitted as permitted by the question.

\begin{center}
\begin{tikzpicture}[circuit logic US, every node/.style={font=\small},
 gate/.style={draw,minimum height=10mm,minimum width=12mm},
 ff/.style={draw,minimum height=13mm,minimum width=15mm}]
\node[and gate US,gate,logic gate inputs=nn] (p1) at (2.5,2.3) {};
\node[and gate US,gate,minimum height=15mm,logic gate inputs=inn] (p2) at (2.5,0) {};
\node[and gate US,gate,minimum height=20mm,logic gate inputs=niii] (p3) at (2.5,-2.3) {};
\foreach \pin/\lab in {1/s_1,2/s_0}
  \draw (p1.input \pin) -- ++(-1.2,0) node[left] {\(\lab\)};
\foreach \pin/\lab in {1/s_1,2/s_0,3/F}
  \draw (p2.input \pin) -- ++(-1.2,0) node[left] {\(\lab\)};
\foreach \pin/\lab in {1/s_1,2/s_0,3/D,4/F}
  \draw (p3.input \pin) -- ++(-1.2,0) node[left] {\(\lab\)};
\node[or gate US,gate,logic gate inputs=nnn] (sum) at (5.6,0) {};
\draw (p1.output) -- ++(.7,0) |- (sum.input 1);
\draw (p2.output) -- (sum.input 2);
\draw (p3.output) -- ++(.7,0) |- (sum.input 3);
\node[and gate US,gate,logic gate inputs=ni] (rst) at (8,0) {};
\draw (sum.output) -- (rst.input 1);
\draw (rst.input 2) -- ++(-.6,0) -- ++(0,-.9) node[below] {\(R\)};
\node[ff] (ff) at (10.5,0) {DFF};
\node[font=\scriptsize,anchor=west] at (ff.west) {D};
\node[font=\scriptsize,anchor=east] at (ff.east) {Q};
\draw (rst.output) -- (ff.west) node[midway,above] {\(s_1^+\)};
\draw (ff.east) -- ++(.8,0) node[right] {\(s_1\)};
\end{tikzpicture}
\end{center}

\begin{center}
\begin{tikzpicture}[circuit logic US,every node/.style={font=\small},
 gate/.style={draw,minimum height=10mm,minimum width=12mm},
 ff/.style={draw,minimum height=13mm,minimum width=15mm}]
\node[or gate US,gate,logic gate inputs=ii] (sum) at (1.5,0) {};
\draw (sum.input 1) -- ++(-.7,0) node[left] {\(s_1\)};
\draw (sum.input 2) -- ++(-.7,0) node[left] {\(T\)};
\node[and gate US,gate,logic gate inputs=nn] (prod) at (4,0) {};
\draw (sum.output) -- (prod.input 1);
\draw (prod.input 2) -- ++(-.5,0) -- ++(0,-.8) node[below] {\(s_0\)};
\node[or gate US,gate,logic gate inputs=nn] (fsum) at (6.2,0) {};
\draw (prod.output) -- (fsum.input 1);
\draw (fsum.input 2) -- ++(-.5,0) -- ++(0,-.8) node[below] {\(D\)};
\node[and gate US,gate,logic gate inputs=ni] (rst) at (8.4,0) {};
\draw (fsum.output) -- (rst.input 1);
\draw (rst.input 2) -- ++(-.5,0) -- ++(0,-.8) node[below] {\(R\)};
\node[ff] (ff) at (10.7,0) {DFF};
\node[font=\scriptsize,anchor=west] at (ff.west) {D};
\node[font=\scriptsize,anchor=east] at (ff.east) {Q};
\draw (rst.output) -- (ff.west) node[midway,above] {\(s_0^+\)};
\draw (ff.east) -- ++(.65,0) node[right] {\(s_0\)};
\end{tikzpicture}
\end{center}

\begin{center}
\begin{tikzpicture}[circuit logic US,every node/.style={font=\small},
 gate/.style={draw,minimum height=10mm,minimum width=12mm}]
\node[and gate US,gate,logic gate inputs=in] (op) at (1,0) {};
\draw (op.input 1) -- ++(-1,0) node[left] {\(s_1\)};
\draw (op.input 2) -- ++(-1,0) node[left] {\(s_0\)};
\draw (op.output) -- ++(1,0) node[right] {\(\mathrm{OP}\)};
\node[and gate US,gate,logic gate inputs=ni] (cl) at (7,0) {};
\draw (cl.input 1) -- ++(-1,0) node[left] {\(s_1\)};
\draw (cl.input 2) -- ++(-1,0) node[left] {\(s_0\)};
\draw (cl.output) -- ++(1,0) node[right] {\(\mathrm{CL}\)};
\end{tikzpicture}
\end{center}
The two state-register circuits have no hidden storage. Their feedback
comes from the labeled \(s_1,s_0\) Q nets.

\clearpage
\section{A five-second timer with no off-by-one error}
A 13-bit unsigned register \(c=c_{12}\cdots c_0\) is sufficient because
\(2^{12}<5000<2^{13}\). During an uninterrupted open interval, count
completed absent-person samples. Define
\[
H=s_1s_0,\qquad
T=[c=4999],\qquad
E=\overline R\,H\,\overline D\,\overline T.
\]
Here \(E\) enables incrementing. When \(E=0\), the counter is
\emph{cleared}, not held:
\[
\boxed{c^+=
\begin{cases}
c+1,&E=1,\\
0,&E=0.
\end{cases}}
\]
Because \(H=0\) before the edge that enters \(\mathsf H\), the timer is
zero at that entry edge. The next 4999 absent samples produce the counts
1 through 4999. At the following, 5000th absent sample,
\(D=0,T=1\), so the state becomes \(\mathsf X\) and the counter clears.
The close command begins exactly \(5000\,\mathrm{ms}=5\,\mathrm{s}\)
after the entry edge. A detection in \(\mathsf H\) resets the same
sequence.

\begin{center}
\begin{tabular}{rllll}
\toprule
Edge after restart & Pre-edge count & \(D\) & Post-edge state & Post-edge count\\
\midrule
0 & --- & --- & \(\mathsf H\) & 0\\
1 & 0 & 0 & \(\mathsf H\) & 1\\
4999 & 4998 & 0 & \(\mathsf H\) & 4999\\
5000 & 4999 & 0 & \(\mathsf X\) & 0\\
\midrule
Any edge in \(\mathsf H\) & Any reachable count & 1 & \(\mathsf H\) & 0\\
\bottomrule
\end{tabular}
\end{center}
Thus \(T=1\) alone does not request closing: the transition also requires
\(D=0\). A person arriving on the timeout edge keeps the door open.
For asynchronous physical arrival times there would be sampling
quantization; the exam explicitly supplies synchronous inputs.

\subsection{Terminal-count logic}
The required bit pattern is
\[
4999=4096+512+256+128+4+2+1=(1001110000111)_2.
\]
The following four AND gates decode every bit. Input bubbles provide the
six required complements.
\begin{center}
\begin{tikzpicture}[circuit logic US,every node/.style={font=\small},
 gate/.style={draw,minimum height=18mm,minimum width=12mm}]
\node[and gate US,gate,logic gate inputs=niin] (g3) at (2,2.4) {};
\node[and gate US,gate,logic gate inputs=nnii] (g2) at (2,0) {};
\node[and gate US,gate,logic gate inputs=iinn] (g1) at (2,-2.4) {};
\foreach \pin/\lab in {1/12,2/11,3/10,4/9}
 \draw (g3.input \pin) -- ++(-1,0) node[left] {\(c_{\lab}\)};
\foreach \pin/\lab in {1/8,2/7,3/6,4/5}
 \draw (g2.input \pin) -- ++(-1,0) node[left] {\(c_{\lab}\)};
\foreach \pin/\lab in {1/4,2/3,3/2,4/1}
 \draw (g1.input \pin) -- ++(-1,0) node[left] {\(c_{\lab}\)};
\node[and gate US,gate,logic gate inputs=nnnn] (term) at (5.8,0) {};
\draw (g3.output) -- ++(.9,0) |- (term.input 1);
\draw (g2.output) -- (term.input 2);
\draw (g1.output) -- ++(1.3,0) |- (term.input 3);
\draw (term.input 4) -- ++(-.45,0) -- ++(0,-1.1) node[below] {\(c_0\)};
\draw (term.output) -- ++(1,0) node[right] {\(T\)};
\end{tikzpicture}
\end{center}

\clearpage
\section{Expanded counter circuit}
An incrementer has carry inputs
\[
p_0=1,\qquad p_1=c_0,\qquad
p_2=c_1c_0,\quad \ldots,\quad
p_{12}=c_{11}c_{10}\cdots c_0.
\]
Equivalently \(p_{i+1}=p_i c_i\) for \(i=0,\ldots,11\).
These are combinational carry signals, not registers. Thus every counter
bit has the explicit next-state equation
\[
\boxed{c_i^+=E(c_i\mathbin{\oplus}p_i),\qquad i=0,\ldots,12.}
\]
The AND with \(E\) implements the synchronous clear described above.
The carry network may be built from the displayed AND recurrence or
treated as the combinational part of an incrementer, as allowed by the
question. All thirteen stored bits and their XOR/clear gates appear
below; none is represented by an unexplained counter block.

\begin{center}
\begin{tikzpicture}[circuit logic US,every node/.style={font=\small},
 gate/.style={draw,minimum height=7mm,minimum width=8mm},
 ff/.style={draw,minimum width=14mm,minimum height=6mm}]
\foreach \i in {0,...,12}{
 \pgfmathsetmacro{\yy}{-0.90*\i}
 \node[xor gate US,gate,logic gate inputs=nn] (xor\i) at (2,\yy) {};
 \draw (xor\i.input 1) -- ++(-.7,0) node[left,font=\scriptsize] {\(c_{\i}\)};
 \draw (xor\i.input 2) -- ++(-.7,0) node[left,font=\scriptsize] {\(p_{\i}\)};
 \node[and gate US,gate,logic gate inputs=nn] (and\i) at (4.3,\yy) {};
 \draw (xor\i.output) -- (and\i.input 1);
 \draw (and\i.input 2) -- ++(-.4,0) node[left,font=\scriptsize] {\(E\)};
 \node[ff] (ff\i) at (6.6,\yy) {DFF};
 \node[font=\tiny,anchor=west] at (ff\i.west) {D};
 \node[font=\tiny,anchor=east] at (ff\i.east) {Q};
 \draw (and\i.output) -- (ff\i.west) node[midway,above,font=\scriptsize] {\(c_{\i}^+\)};
 \draw (ff\i.east) -- ++(.7,0) node[right] {\(c_{\i}\)};
}
\end{tikzpicture}
\end{center}
Every \(c_i\) net on a gate input connects to the corresponding Q output.
All \(E\) nets share the same enable. The next section specifies its
gate, and the labeled feedback connections complete the timer circuit.

\clearpage
\section{Enable gate and complete interconnection}
\begin{center}
\begin{tikzpicture}[circuit logic US,every node/.style={font=\small}]
\node[and gate US,draw,minimum height=23mm,minimum width=15mm,
 logic gate inputs=innii] (en) at (0,0) {};
\foreach \pin/\lab in {1/R,2/s_1,3/s_0,4/D,5/T}
 \draw (en.input \pin) -- ++(-1.4,0) node[left] {\(\lab\)};
\draw (en.output) -- ++(1.4,0) node[right] {\(E\)};
\end{tikzpicture}
\end{center}
This implements \(E=\overline R s_1s_0\overline D\,\overline T\).
The complete interconnection below is a map of the already-expanded
subcircuits, not a substitute for their flip-flop drawings.

\begin{center}
\begin{tikzpicture}[>=Latex,every node/.style={font=\small}]
\node[block,minimum width=38mm,minimum height=22mm] (fsm) at (0,0)
 {State equations\\and 2 DFFs};
\node[block,minimum width=37mm,minimum height=22mm] (count) at (6,0)
 {13 counter bit cells\\and carry logic};
\node[block,minimum width=30mm] (dec) at (6,-3) {4999 decoder};
\node[block,minimum width=30mm] (enable) at (0,-3) {Enable AND gate};
\node[block,minimum width=36mm] (out) at (0,2.5) {Output decoder};
\draw[signal] (-3.5,.5) node[left] {\(D,F,R\)} -- (fsm.west |- 0,.5);
\draw[signal] (fsm.north) -- (out.south) node[midway,right] {\(s_1,s_0\)};
\draw[signal] (out.east) -- ++(2,0) node[right] {\(\mathrm{OP},\mathrm{CL}\)};
\draw[signal] (fsm.south) -- (enable.north) node[midway,left] {\(s_1,s_0\)};
\draw[signal] (-3.5,-3) node[left] {\(R,D\)} -- (enable.west);
\draw[signal] (enable.south) -- (0,-4.4) -- (8.5,-4.4)
 -- (8.5,0) -- (count.east) node[pos=.35,below] {\(E\)};
\draw[signal] (count.south) -- (dec.north) node[midway,right] {\(c_{12:0}\)};
\draw[signal] (dec.west) -- (enable.east) node[midway,below] {\(T\)};
\fill (3,-3) circle (1.2pt);
\draw[signal] (3,-3) |- (fsm.east) node[pos=.7,above] {\(T\)};
\end{tikzpicture}
\end{center}
The state logic receives \(D,F,T,R\). The counter needs \(E\) and its own
Q feedback; the decoder reads all thirteen Q bits. There are exactly
15 flip-flops: two for state and thirteen for time. No separate clock
divider, latch, or hidden timer register is required.

\section{Verification and representative traces}
The accompanying \texttt{tools/verify\_final.py} compares the Boolean
next-state equations with a separately written transition-table model
for every state and every combination of \(D,F,T,R\). It also checks
the XOR/AND counter circuit against arithmetic incrementing for every
13-bit pattern, and checks all reachable timer values with both input
bits. It verifies the ideal synchronous examination circuit without
requiring a logic-simulation application. Selected traces include:
\begin{enumerate}[itemsep=2pt,topsep=3pt]
\item Closed with no detection: remain closed, with zero commands and
zero timer.
\item Detection, opening, endpoint pulse: enter \(\mathsf H\) with
timer zero.
\item 4999 absent samples: remain open. The 5000th absent sample starts
closing.
\item Detection on the would-be timeout edge: remain open and clear the
timer; require another full 5000 absent samples.
\item Detection while closing: enter \(\mathsf U\), with
\((\mathrm{OP},\mathrm{CL})=(1,0)\), including when \(F=1\) simultaneously.
\item Closing endpoint with \(D=0\): return to \(\mathsf C\).
\item Initialization from any state: return to \(\mathsf C\) and clear
all timer bits.
\end{enumerate}
